%%%%%%%%%%%%%%%%%%%%%%%%%%%%%%%%%%%%%%%%%%%%%%%%%%%%%%%%%%%%%%%%%%%%%%%%%%%%%%%
%  Correctness review, round 2 (revised manuscript)
%
%  "Mode Coverage in Normalizing Flow Boltzmann Generators via Log-Ratio
%   Variation"  (Feng, Lai, Qi, Ye)  --  formerly "Forward KL Can Be Better"
%
%  Standalone: compiles with pdflatex. Manuscript objects are referred to by
%  descriptive name; the reference key maps each name to its \label.
%%%%%%%%%%%%%%%%%%%%%%%%%%%%%%%%%%%%%%%%%%%%%%%%%%%%%%%%%%%%%%%%%%%%%%%%%%%%%%%
\documentclass[11pt]{article}

\usepackage[margin=1.05in]{geometry}
\usepackage{amsmath,amssymb,amsthm}
\usepackage{enumitem}
\usepackage{booktabs}
\usepackage{longtable}
\usepackage{xcolor}
\usepackage[colorlinks=true,linkcolor=blue!55!black,urlcolor=blue!55!black]{hyperref}

\newcommand{\R}{\mathbb R}
\newcommand{\E}{\mathbb E}
\newcommand{\D}{\,\mathrm d}
\newcommand{\KL}{\mathrm{KL}}
\newcommand{\ESS}{\mathrm{ESS}}
\newcommand{\X}{\mathcal X}
\newcommand{\osc}{\operatorname{osc}}
\newcommand{\Le}{\leqslant}
\newcommand{\Ge}{\geqslant}
\newcommand{\lab}[1]{\texttt{\small #1}}
\usepackage{pifont}
\newcommand{\ok}{\textcolor{green!50!black}{\ding{51}}}

\newcommand{\where}[1]{\smallskip\noindent\textbf{Where.}\ #1\par}
\newcommand{\what}[1]{\smallskip\noindent\textbf{What.}\ #1\par}
\newcommand{\why}[1]{\smallskip\noindent\textbf{Why it matters.}\ #1\par}
\newcommand{\fix}[1]{\smallskip\noindent\textbf{Suggested fix.}\ #1\par}

\title{\bf Correctness review, round 2\\[3pt]
\large \emph{Mode Coverage in Normalizing Flow Boltzmann Generators via
Log-Ratio Variation} --- revised manuscript}
\author{}
\date{\today}

\begin{document}
\maketitle
\vspace{-3em}

\section*{Summary}

The revision addresses seven of the eight round-1 issues and two of the round-1
minor items, and it does so well: the one-nat cap is now in the main text, the
unbiasedness hypothesis travels with every statement of the result, the
inference pass is Algorithm~6 with an explicit scope remark, the
$\hat F_\Sigma/(2\sqrt N)$ scale is applied to all three chiral targets, the
$\tau_{\mathrm v}$ default matches the runs, and the framing of \S2.2 is now
honest about whose rate it is. The mathematics was sound before and is
unchanged. Three things in the \emph{new} material need attention before
submission, in this order:

\begin{enumerate}[leftmargin=*,itemsep=2pt]
\item the AI-use declaration became \emph{narrower} in a revision in which the
      AI-assisted content became \emph{broader}, and it still lacks the sentence
      SIAM requires verbatim (\S\ref{sec:decl});
\item the new ``time average'' paragraph after Theorem~2.1 is a heuristic that
      does not hold; the correct reconciliation is shorter and is given below
      (\S\ref{sec:timeavg});
\item the new paragraphs quoting numbers from the FAB and TA-BG papers are
      almost all right --- I checked them against the sources --- but two
      details need adjusting (\S\ref{sec:cites}).
\end{enumerate}

\noindent
The rest is editorial: stale sentences left over from the reframing, one
insertion that broke the flow of a remark, and the round-1 minor items that
were not picked up.

\subsection*{Round-1 status}

\begin{center}
\small
\begin{tabular}{@{}clll@{}}
\toprule
\# & round-1 item & status & note \\
\midrule
1 & dangling reference; one-nat cap unstated & \ok\ fixed & \S2.3, blue paragraph \\
2 & unbiasedness without hypothesis & \ok\ fixed & abstract, contribution 4, \S7 \\
3 & Remark B.1 excludes defaults & \ok\ fixed & but see \S\ref{sec:remB1}: the insertion breaks the remark \\
4 & $e^{-t}$ rate is forward KL's & \ok\ fixed & but \S2 intro and App.~A title are stale (\S\ref{sec:stale}) \\
5 & finite-sample story, one target only & \ok\ fixed & option (ii), all three scales \\
6 & $q$ double-booked & \ok\ fixed & but $Q_i$ now meets $Q_k$ (\S\ref{sec:Q}) \\
7 & theorem attached to unnumbered scheme & \ok\ fixed & Algorithm~6 + scope remark \\
8 & $\tau_{\mathrm v}$ default never used & \ok\ fixed & $0.4$ everywhere \\
\midrule
m1 & achiral ratios round the other way & \ok\ fixed & $2.17$, $1.87$ \\
m2 & leftover commented block & \ok\ fixed & \\
m3 & ``stage time'' vs ``wall time'' & open & Table~10 caption still differs \\
m4 & ``zero'' should be ``constant'' & open & \S1, unchanged \\
m5 & ``decays as'' is a bound & \ok\ fixed & conclusions \\
m6 & tie rule missing from Algorithm~5 & open & \\
m7 & termination not argued & open & \\
m8 & \S6.5 cross-reference to \S6.3 & open & \\
m9 & $E_{\mathrm{vac}}$ defined then dropped & open & \\
m10 & symbol reuse & open & $q_k$ in Lemma C.1 still there \\
\bottomrule
\end{tabular}
\end{center}

\subsection*{Reference key}

\begin{center}
\small
\begin{tabular}{@{}ll@{}}
\toprule
name used in this note & label in the manuscript \\
\midrule
the accuracy theorem & \lab{thm: accuracy} \\
the propagation theorem & \lab{thm: propagation} \\
the domination remark & \lab{rem: appendix-domination} \\
the one-step lemma & \lab{lem: onestep} \\
the FR flow & \lab{eq: FR-flow} \\
the dissipation identity & \lab{eq: combined-rate} \\
the stationarity identity & \lab{eq: surrogate-stationary} \\
the biased loss & \lab{eq: KLXX-biased} \\
the training algorithm & \lab{alg: training} \\
the SMC surrogate & \lab{alg: annealing} \\
the staging algorithm & \lab{alg: boltzmann} \\
the inference algorithm & \lab{alg: inference} \\
the FR appendix & \lab{sec: appendix-FR} \\
the discussion section & \lab{subsec: mol-discussion} \\
\bottomrule
\end{tabular}
\end{center}

%==============================================================================
\section{The declaration of AI use}
\label{sec:decl}
%==============================================================================

\where{``Declaration of AI use'', after Data Availability.}

\what{The previous version declared assistance with ``writing and testing the
code, developing the proof of Theorem~[propagation], and general manuscript
writing.'' The revision declares assistance with ``writing and testing the
codes'' only. Between the two versions, several passages that an AI drafted or
reviewed entered the manuscript --- the positivity argument in the domination
remark, the scope remark after the propagation theorem, the hypothesis clause
on unbiasedness, the three-target scale estimates in \S6.5, and the corrected
ratios --- and the whole manuscript was reviewed twice. The declared scope
moved in the opposite direction from the actual one.}

\why{Under the SIAM policy you asked about, disclosure is mandatory, the
statement ``The authors assume responsibility for all content'' is required
verbatim, and ``negligent or undisclosed usage of AI, whether intentional or
not, is grounds for an author integrity investigation.'' Items~2.2 (proofs) and
2.1 (related-work accuracy) are called out specifically, and the revision
touches both: the propagation theorem's exposition and the new \S1.2 and \S6.6
paragraphs that quote other papers' numbers. Nothing here is a problem
scientifically --- the policy permits all of it --- but only if it is declared.
The current sentence ``take full responsibility for the correctness of the
results'' is narrower than the required one (results, not content) and the
narrowed scope is now inaccurate. This is the one item in the paper that could
hurt you independently of the science.}

\fix{Restore the fuller scope and add the required sentence. Name whichever
models were actually used; I do not know which ones were used for the code.}

\begin{verbatim}
\subsection*{Declaration of AI use}
The authors used Claude (Anthropic) to assist with writing and
testing the code, with drafting and checking proofs, and with
reviewing and editing the manuscript, including the checking of
numerical claims quoted from other publications. All data and
numerical results reported in this paper are generated directly by
our Python scripts and are presented without any embellishment. The
authors have reviewed all AI-assisted content. The authors assume
responsibility for all content.
\end{verbatim}

%==============================================================================
\section{The ``time average'' paragraph after Theorem 2.1}
\label{sec:timeavg}
%==============================================================================

\where{\S2.3, the last blue paragraph before \S3.}

\what{The paragraph reconciles the fixed $\tilde\pi$ of the accuracy theorem
with the per-step surrogate of the training algorithm by saying that ``the
surrogate the flow descends against over a training run is the time average of
the per-step SMC sample sets, and it is that averaged distribution that plays
the role of $\tilde\pi$ in the theorem.''}

\why{This does not hold, for three reasons. (a) The per-step surrogate is built
from pushforward samples of the \emph{current} flow, so its time average along
a training run is a trajectory-dependent object, not a distribution the theorem
can be handed. (b) The target variation is quadratic in its outer measure
($\tilde\pi\otimes\tilde\pi$), so the average of the per-step losses is not the
loss at the averaged surrogate; averaging commutes with the KL term only.
(c) The theorem is a statement about a \emph{stationary point}, where no
averaging is needed at all. The correct reconciliation is shorter and is
already implicit in the proof: nothing in Appendix~B uses any property of
$\tilde\pi$ beyond the stationarity identity, exactly as nothing uses any
property of $\xi$ beyond it --- the same observation the paper already makes
for $\bar\nu$.}

\fix{Replace the paragraph by:}

\begin{verbatim}
Theorem~\ref{thm: accuracy} holds a single $\tilde\pi$ fixed, while
Algorithm~\ref{alg: training} rebuilds the surrogate at every step
from the current flow. The two meet at a stationary point. The SMC
samples enter the loss as data and are never differentiated through,
so the gradient that vanishes at a stationary flow $\nu_\star$ is the
Fisher--Rao gradient of \eqref{eq: KLXX-biased} with $\tilde\pi$
frozen at the surrogate Algorithm~\ref{alg: annealing} produces when
initialized from $\nu_\star$, namely the marginal law $\tilde\pi_\star$
of one output particle. Identity \eqref{eq: surrogate-stationary}
therefore holds with $\tilde\pi=\tilde\pi_\star$, and the proof of
Theorem~\ref{thm: accuracy} uses nothing about $\tilde\pi$ beyond that
identity, so the theorem applies with that surrogate, just as it
applies with $\xi$ containing the detached copy of $\nu_\star$.
\end{verbatim}

\noindent
One further honesty note, optional: the empirical variation pairs two particles
of the same SMC batch, which after resampling are exchangeable but not
independent, so $\tilde\pi_\star\otimes\tilde\pi_\star$ is the independent-pair
idealization of that batch. This was equally true of the previous version and
is worth one clause, not a paragraph.

%==============================================================================
\section{Claims quoted from other papers}
\label{sec:cites}
%==============================================================================

The new \S1.2 paragraph and \S6.6 quote specific numbers from
the FAB paper (arXiv:2208.01893, Table~2) and TA-BG (arXiv:2501.19077, Table~1).
I checked each against the source. Most are exact; two need adjusting and one
is confirmed for only one of the two papers it is attributed to.

\begin{center}
\small
\setlength{\tabcolsep}{4pt}
\begin{longtable}{@{}p{0.43\textwidth}p{0.36\textwidth}p{0.13\textwidth}@{}}
\toprule
claim in the manuscript & source & verdict \\
\midrule
\endhead
FAB w/ buffer ESS $92.8\%$ on alanine dipeptide (FAB Table~2) & $92.8\pm0.1$ & \ok \\
reverse KL flow $54\%$ vs.\ ML flow $2.8\%$ (FAB Table~2) & $54\pm12$ vs.\ $2.8\pm0.6$ & \ok \\
reverse KL Ramachandran KLD worse ``by more than two orders of magnitude'' & $3.17\pm0.20$ vs.\ $(3.42\pm0.45)\times10^{-3}$ (FAB w/ buffer), $(7.57\pm3.80)\times10^{-3}$ (ML) & \ok \\
FAB paper calls the ESS of mode-missing baselines ``spurious'' & ``\dots and spurious values for the ESS'' & \ok\ (verbatim) \\
FAB retrained: $94.81\%$ (TA-BG Table~1) & $(94.81\pm0.04)\%$ & \ok \\
TA-BG ``reports the ESS on flow samples even though that quantity does not capture mode collapse'' & ``we chose to only use the reverse ESS, even though it does not capture mode collapse'' & \ok \\
TA-BG Table~1: reverse KL has the highest ESS on tetrapeptide and hexapeptide & tetra: $74.88$ vs.\ $63.59$ (FAB), $62.36$ (TA-BG), $45.29$ (fwd); hexa: $22.22$ vs.\ $14.71$, $14.64$, $10.97$ & \ok \\
\dots ``while collapsing almost entirely onto the main mode'' & TA-BG says ``now almost entirely collapses'' of the \emph{hexapeptide}; for the tetrapeptide it gives RAM KLD $0.30$ vs.\ $\sim\!0.003$--$0.007$ & adjust wording \\
FAB ``resolving the metastable states of its largest system only partially'' & ``While FAB is partially able to resolve the metastable states, they differ in shape'' & \ok \\
FAB ``reaching comparable accuracy on peptides at roughly three times the target energy evaluations'' & PE evals FAB/TA-BG: dipeptide $2.13{\times}10^8/7.56{\times}10^7=2.8$; tetrapeptide $2.8$; hexapeptide $4.2{\times}10^8/3.08{\times}10^8=1.4$ & adjust: holds for di- and tetrapeptide only \\
``both computed \dots after clipping the highest $0.01\%$ of the importance weights'' & TA-BG: ``we clipped the 0.01\,\% highest importance weights to the lowest value among them'' \ok. FAB paper: no clipping statement found in two passes over the PDF & confirm for FAB, or attribute to TA-BG only \\
\bottomrule
\end{longtable}
\end{center}

\fix{Three edits. (i) \S1.2: ``\dots on alanine di- and tetrapeptide at
roughly three times the target energy evaluations, and resolving the metastable
states of its largest system only partially.'' (ii) \S6.6: ``\dots the reverse
KL run has the highest ESS of the four methods on both alanine tetrapeptide and
alanine hexapeptide, with a Ramachandran KL divergence two orders of magnitude
worse than the others' on both and, on the hexapeptide, an almost complete
collapse onto the main mode.'' (iii) \S6.6, clipping: either locate the
statement in the FAB paper and keep ``both'', or write ``the latter computed
after clipping the highest $0.01\%$ of the importance weights''. Under
SIAM~2.1 the authors are responsible for these being right, so this is worth
five minutes with the two PDFs.}

%==============================================================================
\section{Remark B.1: the insertion broke the remark}
\label{sec:remB1}
%==============================================================================

\where{The domination remark, Appendix~B.}

\what{The purple positivity argument was inserted between ``that is
$\nu_\star\Ge(1-2\lambda-2c\vartheta)\tilde\pi>0$'' and ``Dividing $\pi$ by
this bound and integrating gives \dots''. After the insertion, ``this bound''
refers back across five lines that have just said the bound is not needed, and
the two conclusions $\E_{\tilde\pi}[w_\star]\Le(1-2\lambda-2c\vartheta)^{-1}$,
$\E_\xi[w_\star]\Le c(\cdots)^{-1}$ now read as consequences of the
self-enforcing positivity, which they are not --- they need the quantitative
domination bound.}

\fix{Move the purple sentences to the end of the remark, after the two
expectation bounds, and open them with ``The condition is sufficient but
conservative: \dots''. Consider also adding one clause after the FR flow in
\S2.2, since the flow case is where a reader first meets the coefficient
$1+2\lambda S^\pi$ and wonders about its sign: ``where $\nu_t$ is small the
weight is large, both correlations approach $+1$, and the drift approaches
$\pi(1+2\lambda)+2\vartheta\xi>0$, so positivity is preserved at every
$\lambda,\vartheta\Ge0$.''}

%==============================================================================
\section{Stale sentences from the reframing of \S2.2}
\label{sec:stale}
%==============================================================================

\S2.2 was retitled ``Fisher--Rao gradient of the KLXX loss'' and its opening
now asks ``what the gradient flow it induces dissipates''. Three places did not
follow:

\begin{enumerate}[leftmargin=*,itemsep=2pt]
\item \S2, the paragraph before \S2.1: ``The three subsections below take the
      generalized loss in turn: a formal account of how the mixture variation
      contributes to mode discovery, \emph{a derivation of exponential
      convergence for the Fisher--Rao gradient flow}, and an accuracy bound
      \dots'' --- replace the middle clause by ``the Fisher--Rao gradient flow
      and its dissipation identity''.
\item The FR appendix: its title ``Proof of Fisher--Rao gradient flow
      convergence'' and its first paragraph (``ask how fast $\nu_t$ approaches
      the target'') still advertise the old framing. Retitle ``Fisher--Rao
      gradient flow and dissipation identity''; the derivation itself can stay
      as it is, including the closing $e^{-t}$ sentence, which is now
      correctly qualified.
\item \S2.2 has the detachment of $\bar\nu$ twice in a row: the new blue
      sentence ``Since $\bar\nu$ is detached, $\xi$ is held fixed when the
      gradient is taken and then follows $\nu_t$ along the flow \dots'' is
      immediately followed by the old ``The mixture $\xi$ depends on $\nu_t$
      through $\bar\nu$, but $\bar\nu$ is detached and takes no part in the
      gradient calculation, so $\xi$ is held fixed when the variation is
      taken.'' Delete the second.
\end{enumerate}

\noindent
Also in \S2.2, in the blue text after the dissipation identity: ``At
$\lambda=\vartheta=0$ the flow is the linear equation \dots and the dissipation
is the single term $-\chi^2$. Since $\chi^2\Ge\KL$, \emph{discarding them}
leaves the decay \dots'' --- the antecedent of ``them'' is two sentences back;
write ``discarding the two variation terms''.

%==============================================================================
\section{Other new material}
%==============================================================================

\subsection{\S5.1, the FAB comparison}

\begin{itemize}[leftmargin=*,itemsep=2pt]
\item ``at convergence FAB should hold the highest ESS here, and on Three-Well
      it does, at $0.985$'' --- $\KL$+$\X_\pi$ has $0.994$ on Three-Well. The
      table's bold rule (highest among mode-complete methods) is right; the
      sentence needs the same qualifier: ``the highest ESS among the
      mode-complete methods, and on Three-Well it does''.
\item ``we initialize the same neural spline flow at the identity and compare
      the four losses above'' --- five methods now; ``the four losses above and
      FAB''.
\item ``so the runs differ only in the objective'' --- FAB also brings its AIS
      path to $\pi^2/\nu$ and the prioritized replay buffer, which the
      sentence before lists. ``differ only in the objective and the sampling
      machinery that objective requires'' is accurate.
\item The four figure captions now begin with FAB; the Two-Moon paragraph's
      ``all four losses attain coverage $1$'' is still true of the KL family
      and FAB alike, so it can stand.
\end{itemize}

\subsection{\S1.2, the new comparison paragraph}

\begin{itemize}[leftmargin=*,itemsep=2pt]
\item ``LDR-L1 is within a factor of two of $\KL$+$\X_\pi$ at the same
      coefficient, so it is already represented among the losses we run'' ---
      a two-sided bound on the \emph{penalty} does not make one loss represent
      the other in training, and the paper said as much a page earlier (``This
      suggests similar training performance, although the losses are not
      equivalent''). ``so it is approximately represented'' keeps both
      sentences consistent.
\item ``rather than against the methods of Section~[loss-related]'' --- the
      paragraph is inside that section. ``the other methods of this section''.
\end{itemize}

\subsection{\S2.2 remark}

``Fisher--Rao geometry works in unconstrained density space and therefore
cannot see mode collapse, which is a pathology of the parametrization and of
the optimizer.'' --- It is at least as much a pathology of the \emph{sampled
surrogate}: the exact forward KL in density space sees every mode, and it is
the missing training points that lose them. Theorem~2.1 is precisely the
density-space model of that part. ``a pathology of the sampled surrogate, the
parametrization, and the optimizer''.

\subsection{\S4.4, the Feynman--Kac paragraph}

``with the trained stage map composed with the rejuvenation kernel as the
mutation and the stage weight $w_k$ as the potential'' --- in the standard
mutation-then-selection form, the mutation that precedes selection by $w_k$ is
$Q_{k-1}$ followed by $G_k^{-1}$ (the previous stage's kernel, then this
stage's map), since $w_k$ is evaluated at $Y_k=G_k^{-1}(X_{k-1})$. One
clause: ``the rejuvenation kernel of one stage composed with the map of the
next as the mutation''.

\subsection{$Q_i$ and $Q_k$}
\label{sec:Q}

The round-1 suggestion to rename the charges $Q_i$ assumed the Markov kernel
$Q_k$ stayed in Appendix~C. Algorithm~6 has now, rightly, moved $Q_k$ into
\S4.4, so the two meet in the main text after all --- my suggestion, my
apology. Use $c_i$ or $\mathsf q_i$ for the charges. And $q_k$ in the proof of
the one-step lemma is still there; $\delta_k$.

%==============================================================================
\section{Minor}
%==============================================================================

\begin{enumerate}[leftmargin=*,itemsep=4pt]
\item \textbf{Empty section.} \verb|\section*{Acknowledgments}| now has no
      body; either fill it or delete it.
\item \textbf{Round-1 minor items still open:} ``almost surely \emph{zero} on
      the support of $\omega$'' (\S1) should be \emph{constant}; Table~10 says
      ``total stage time'' while Tables~8--9 say ``total wall time'' and \S1.1
      quotes $72.7$ minutes as a completion time; the tie rule (``ties
      retaining the identity'', \S6.2) is still absent from Algorithm~5; the
      termination of Algorithm~5 is still unargued; \S6.5 still points at
      \S6.3 for the chiral references; $E_{\mathrm{vac}}$ is still defined and
      then dropped; $t_{\mathrm{safe}}=0.2$ is still a default no run uses.
\item \textbf{Wording.} \S1: ``A more universal strategy'' --- the previous
      ``A complementary strategy'' was better English. \S2: ``with no further
      confusion'' --- ``without further comment''.
\item \textbf{Punctuation.} In Appendix~C the display of the recurrence ends
      ``with the second term absent at $k=0$,'' and the next sentence begins
      ``At $k=0$''; the comma should be a period. Appendix~C's
      \verb|\section{ Proof of ...}| has a leading space left from the removed
      color command.
\item \textbf{Before submission.} Strip the blue/purple markup; hyperref will
      also warn about \verb|\color| inside the \S2.2 \verb|\subsection| title
      and the Algorithm~6 \verb|\caption| (moving arguments), harmlessly, until
      it is removed.
\end{enumerate}

%==============================================================================
\appendix
\section{Numbers re-checked in this round}
%==============================================================================

\begin{center}
\small
\setlength{\tabcolsep}{4pt}
\begin{tabular}{@{}p{0.50\textwidth}p{0.20\textwidth}p{0.20\textwidth}@{}}
\toprule
quantity & in paper & recomputed \\
\midrule
achiral wall-time ratios & $2.02,2.17,1.87$ & $2.017,2.170,1.866$ \\
$\hat F_\Sigma/(2\sqrt N)$, ADP / butanediol / Ac-Pro & $0.011/0.006/0.023$ & $0.0108/0.0059/0.0228$ \\
FAB coverage as a mode count: Himmelblau $0.736$, Sparse $0.768$ & 3 of 4 wells & consistent with $3/4$ plus ball overlap \\
FAB Three-Well $0.985$ vs.\ mode-complete KLXX $0.982$, $\hat\pi$-only $0.968$ & FAB bold & consistent with the caption's rule \\
``$0.83$ or above for every other method'' on Two-Moon & --- & min is $0.831$ \\
FAB geometric path from $\nu$ to $\pi^2/\nu$ passes through $\pi$ & stated & $\nu^{1/2}(\pi^2/\nu)^{1/2}=\pi$ \\
$k=0$ term of $F_\Sigma$ equals $\prod_{j=1}^K\ESS(w_j)^{-1/2}$ & stated & confirmed \\
\bottomrule
\end{tabular}
\end{center}

\noindent
Everything verified in round 1 (all propagation factors, stage schedules, atom
counts, $\Delta G$ from populations, and so on) is unchanged in this version
and was not re-derived.

\end{document}
